\everymath{\displaystyle}

%%%%%%%%%%%%%%%%%%%%%%%%%%%%%%%%%%%%%%%%
%           main body
%%%%%%%%%%%%%%%%%%%%%%%%%%%%%%%%%%%%%%%%

%-----------------------------------
%	TITLE PAGE
%-----------------------------------

\title[第五部分\\
循环模, 单模, 有限生成模]{\texorpdfstring{\LARGE \bfseries 第五部分\quad 循环模, 单模, 有限生成模}{第五部分 循环模, 单模, 有限生成模}}
\author[抽象代数 2]{\Large \bfseries 抽象代数 2}
\date{}

%------------------------------------------------

\begin{document}

%------------------------------------------------

\begin{frame}

\titlepage % Print the title page as the first slide

\end{frame}

%------------------------------------------------

\section[生成子模]{生成的子模}

%------------------------------------------------

\begin{frame}{生成的子模}

\begin{Def}[生成的子模]
$R$ 为幺环, $M$ 为一个 $R$-模, $S \subset M$, 令 $\langle S \rangle$
为 $M$ 中所有包含 $S$ 的子模的交, 即
\[
\langle S \rangle = \bigcap_{\substack{S \subset H \subset M \\ H \text{ 为子模}}} H ,
\]
称为 $S$ 在 $M$ 中生成的子模, 且
\[
\langle S \rangle = \Big\{ \sum_{i=1}^{\infty} a_i \cdot \alpha_i \;\Big|\; a_i \in R,\ \alpha_i \in S \Big\}.
\]
\end{Def}

\pause

\begin{itemize}
\item 特别地, 若 $S = \{ \alpha_1, \alpha_2, \dots, \alpha_n \}$ 为 $n$ 个元素的有限集合,
  则 $\langle S \rangle = \langle \alpha_1, \dots, \alpha_n \rangle$,
  称为 $M$ 的一个\alert{有限生成子模};
\item 又若 $S = \{ \alpha \}$, $\langle S \rangle = \langle \alpha \rangle
  = R\alpha = \{ a \cdot \alpha \mid a \in R \}$,
  称为 $M$ 的一个\alert{循环子模}.
\end{itemize}

\end{frame}

%------------------------------------------------

\section[单模]{单模与 Schur 引理}

%------------------------------------------------

\begin{frame}{单模}

\begin{Def}[单模]
若 $M$ 为一个 $R$-模, $M \neq 0$, 且仅有平凡的子模, 即 $H \subset M$ 为子模,
则有 $H = 0$ 或 $M$, 称 $M$ 为 $R$-单模.
\end{Def}

\pause

\begin{itemize}
\item $M$ 为 $R$-单模 $\iff$ $M$ 为 $R$-循环模, 且任一 $M$ 的非零元为生成元.
\end{itemize}

\end{frame}

%------------------------------------------------

\begin{frame}{Schur's 引理}

\begin{lem}[Schur's 引理]
若 $M$ 为 $R$-单模 $\Longrightarrow$ $\operatorname{End}_R(M)$ 为一个域.
\end{lem}

\startproof[bg=yellow, fg=red](证明)
$\forall\ \sigma \in \operatorname{End}_R(M)$, $\sigma(M) \subset M$ 为 $M$
的一个子模, $\because M$ 为单模, 那么若 $\sigma \neq 0$, 则有 $\sigma(M) = M$,
即 $\sigma$ 为满的.
同理, $\ker \sigma \subset M$ 为子模, 由于 $\sigma \neq 0$,
知 $\ker \sigma \neq M$, 所以 $\ker \sigma = 0$, 即 $\sigma$ 是单的.
由此即知 $\sigma$ 为 $M \to M$ 的模同构.
\qed

\pause

\begin{itemize}
\item 注: Schur's 引理的逆命题不真.
\end{itemize}

\end{frame}

%------------------------------------------------

\section[例子]{例子}

%------------------------------------------------

\begin{frame}{理想与子模的积}

\begin{eg}[理想与子模的积]
$M$ 为 $R$-模, $N \subset R$ 为左理想, $H \subset M$ 为一个子模,
令
\[
N \cdot H = \Big\{ \sum_{i=1}^{\infty} a_i \cdot \alpha_i \;\Big|\; a_i \in N,\ \alpha_i \in H \Big\}
\;\Longrightarrow\; N \cdot H \subset M \text{ 为子模},
\]
又称为 $N$ 与 $H$ 所生成的子模.
\end{eg}

\end{frame}

%------------------------------------------------

\begin{frame}{零化子}

\begin{eg}[零化子]
$M$ 为 $R$-模, $\alpha \in M$, 令
$\operatorname{Ann}_R(\alpha) = \{ a \mid a \in R, \text{且 } a \cdot \alpha = 0 \}$,
$\operatorname{Ann}_R(\alpha)$ 为 $R$ 的一个左理想, 又称为 $\alpha$ 的零化子.
且
\[
\operatorname{Ann}_R(M)
= \bigcap_{\substack{\alpha \in M \\ \alpha \neq 0}} \operatorname{Ann}_R(\alpha).
\]
\end{eg}

\end{frame}

%------------------------------------------------

\begin{frame}{循环模的刻画}

\begin{prop}[循环模的刻画]
$R$ 为 $R$-模, $\operatorname{Ann}_R(\alpha)$ 为左理想
$\Longrightarrow$ $R$ 的商模 $R/\operatorname{Ann}_R(\alpha) \cong \langle \alpha \rangle$
($R$-模同构).
\end{prop}

\startproof[bg=yellow, fg=red](证明)
令 $R \xrightarrow{\ \sigma\ } \langle \alpha \rangle$ by
$\sigma(a) = a \cdot \alpha$, 则 $\sigma$ 为 $R$-模同态,
因为
\[
\sigma(a+b) = (a+b)\alpha = a\cdot\alpha + b\alpha = \sigma(a) + \sigma(b),
\]
\[
\sigma(\lambda a) = (\lambda a)\alpha = \lambda(a\alpha) = \lambda\,\sigma(a),
\quad \lambda \in R;
\]
且 $\sigma$ 为满同态.
\pause
$\Longrightarrow$ $\ker \sigma$ 为子模, 且 $R/\ker\sigma \cong \langle \alpha \rangle$,
$\ker \sigma = \operatorname{Ann}_R(\alpha)$.
\qed

\end{frame}

%------------------------------------------------

\begin{frame}{循环模的例子}

\begin{eg}[循环模的例子]
$R$ 为 $R$-模, 则 $R = \langle 1 \rangle$ 为 $R$-循环模;
若 $N \subset R$ 为左理想, $R/N = \langle \bar{1} \rangle$ 为 $R$-循环模.
\end{eg}

\end{frame}

%------------------------------------------------

\begin{frame}{商环 $R/N$ 的性质}

\begin{eg}[商环 $R/N$ 的性质]
$R$ 为幺环, $N \subset R$ 为一个理想, $R/N$ 为商模, 也是幺环,
\begin{itemize}
\item[(i)] $R/N$ 为单模 $\Longrightarrow$ $N$ 为 $R$ 的极大理想;
\item[(ii)] $\operatorname{Ann}_R(R/N) = N$;
\item[(iii)] $\operatorname{End}_R(R/N)$ 反同构于 $R/N$.
\end{itemize}
\end{eg}

\end{frame}

%------------------------------------------------

\begin{frame}{内含理想与理想化子}

\startproof[bg=yellow, fg=red](思路)
若 $I \subset R$ 为一个左理想, 令
\[
(I:R) = \{ a \in R \mid a \cdot R \subset I \}
\quad \text{为 } R \text{ 的理想},
\]
\[
B = \{ a \in R \mid I \cdot a \subset I \}
\quad \text{为 } R \text{ 的子环}.
\]
可以证明, $(I:R) \subset I$, 称为 $I$ 的内含的理想, 且为 $I$ 中内含的一个最大的理想;
$I \subset B \subset R$, $B$ 为子环, 且 $I$ 为 $B$ 中的理想, 称 $B$ 为 $I$
的理想化子, 并且若 $I \subset R_1 \subset R$, $R_1$ 为子环, $I$ 为 $R_1$
的理想, 则必有 $R_1 \subset B$.
\qed

\end{frame}

%------------------------------------------------

\begin{frame}{商模 $R/I$ 的性质}

\begin{eg}[商模 $R/I$ 的性质]
$R$ 为一个幺环, $I \subset R$ 为左理想, 商模 $R/I$ 有如下性质:
\begin{itemize}
\item[(i)] $R/I$ 为 $R$-单模 $\iff$ $I$ 为 $R$ 的极大左理想;
\item[(ii)] $\operatorname{Ann}_R(R/I) = (I:R)$;
\item[(iii)] $\operatorname{End}_R(R/I)$ 反同构于 $B/I$.
\end{itemize}
\end{eg}

\end{frame}

%------------------------------------------------

\begin{frame}{商模 $R/I$ 性质的证明 (i) (ii)}

\startproof[bg=yellow, fg=red]((i) 的证明)
$R \xrightarrow{\ \pi\ } R/I$ 为自然同态 (作为 $R$-模), $\ker \pi = I$,
由同态定理, $R/I$ 为单模 $\iff$ $I$ 为极大左理想.
\qed

\pause

\startproof[bg=yellow, fg=red]((ii) 的证明)
$a \in \operatorname{Ann}_R(R/I) \iff \forall\ \bar{\alpha} \in R/I$,
有 $a \cdot \bar{\alpha} = 0$, 即 $\overline{a \cdot \alpha} = 0$
$\iff$ $a \cdot \alpha \in I$, $\forall\ \alpha \in R$,
即 $a \cdot R \subset I$,
即 $\operatorname{Ann}_R(R/I) = (I:R)$.
\qed

\end{frame}

%------------------------------------------------

\begin{frame}{商模 $R/I$ 性质的证明 (iii)}

\startproof[bg=yellow, fg=red]((iii) 的证明)
$\sigma \in \operatorname{End}_R(R/I)$, $R/I \xrightarrow{\ \sigma\ } R/I$,
又有 $R/I = \langle \bar{1} \rangle$ 为 $R$-循环模.
$\forall\ \bar{\alpha} \in R/I$,
\[
\sigma(\bar{\alpha}) = \sigma(a \cdot \bar{1}) = a\,\sigma(\bar{1}),
\]
即 $\sigma$ 由 $\sigma(\bar{1})$ 所完全决定, 令 $\sigma(\bar{1}) = \bar{b}
\in R/I$, 记 $\sigma$ 为 $\sigma_b$, 其中 $b \in R$;
\pause
反过来, $\forall\ b \in R$, 若能决定 $\sigma_b$,
那么 $\forall\ a \in I \subset R$, $\sigma_b(\bar{\alpha}) = 0$,
但是 $\sigma_b(\bar{\alpha}) = a\,\sigma_b(\bar{1}) = a\bar{b}
= \overline{ab} \Longleftrightarrow ab \in I$,
即 $I \cdot b \subset I$, 即 $b \in B$.
所以 $\operatorname{End}_R(R/I) = \{ \sigma_b \mid b \in B/I \}$.
\pause
下面只需验证 环 $B/I \xrightarrow{\ \varphi\ } \operatorname{End}_R(R/I)$
by $\varphi(\bar{b}) = \sigma_b$, $\varphi$ 为 $1$-$1$ 对应,
且为环的反同态即可.
\qed

\end{frame}

%------------------------------------------------

\section[幺环的根]{幺环的根}

%------------------------------------------------

\begin{frame}{幺环的根}

\begin{Def}[幺环的根]
一个幺环的根 $J(R)$ 或 $\operatorname{rad}(R)$,
\[
J(R) := \bigcap_{I \text{ 为 } R \text{ 的极大左理想}} I ,
\]
$J(R)$ 为 $R$ 的理想.
\end{Def}

\end{frame}

%------------------------------------------------

\begin{frame}{单模与极大左理想}

\begin{lem}[单模与极大左理想]
$R$ 为幺环, 则 $M$ 为 $R$-单模 $\iff$ $\exists\ R$ 的一个极大左理想 $I$,
使得 $R/I \cong M$.
\end{lem}

\startproof[bg=yellow, fg=red](证明)
若 $M \cong R/I$, $I$ 为极大左理想, 则 $R/I$ 为单模, 即 $M$ 也为单模.
反过来, 若 $M$ 为单模, $M = \langle \alpha \rangle$, $\alpha \in M$,
$\alpha \neq 0$, 那么 $M \cong R/\operatorname{Ann}_R(\alpha)$,
$\operatorname{Ann}_R(\alpha)$ 为 $R$ 的极大左理想.
\qed

\end{frame}

%------------------------------------------------

\begin{frame}{单模与极大左理想的对应}

\begin{itemize}
\item $R$ 的单模与极大左理想是 $1$-$1$ 对应的.
\item ① $I \subset R$ 为极大左理想 $\Longrightarrow$ $R/I$ 为单模;
\item ② 若 $M$ 为 $R$-单模, $\alpha \in M$, $\alpha \neq 0$,
  则 $\operatorname{Ann}_R(\alpha)$ 为极大左理想,
  且 $M \cong R/\operatorname{Ann}_R(\alpha)$.
\end{itemize}

\pause

\[
\big\{ R/I \;\big|\; I \subset R \text{ 为极大左理想} \big\}
\text{ 为全部的 } R\text{-单模}.
\]

\end{frame}

%------------------------------------------------

\begin{frame}{根的刻画}

\begin{lem}[根的刻画]
$R$ 为一个幺环, 则
\[
J(R) = \bigcap_{I \text{ 为 } R \text{ 的极大左理想}} (I:R) ,
\]
即知 $J(R)$ 为理想.
\end{lem}

\startproof[bg=yellow, fg=red](证明)
由定义即知
$\bigcap\limits_{I \text{ 为 } R \text{ 的极大左理想}} (I:R) \subset J(R)$,
因为 $(I:R) \subset I$.
\pause
反过来,
\[
\bigcap_{I \text{ 为极大左理想}} (I:R)
= \bigcap_{I \text{ 为极大左理想}} \operatorname{Ann}_R(R/I)
= \bigcap_{M \text{ 为 } R \text{ 单模}} \operatorname{Ann}_R(M)
\]
\[
= \bigcap_{M \text{ 为 } R \text{ 单模}}\ \bigcap_{\substack{\alpha \in M \\ \alpha \neq 0}}
\operatorname{Ann}_R(\alpha)
\ \supset\ J(R)
\]
因为每一个 $\operatorname{Ann}_R(\alpha)$ 都是 $R$ 的一个极大左理想.
\qed

\end{frame}

%------------------------------------------------

\end{document}
